\section{Make a local coefficient fibre the foundational theorem}
\label{sec:coefficients}

\subsection{The mathematical statement}

Let \(I\) be a finite index type, \(R\) a commutative semiring, and
\(w:I\to\N\), \(u:I\to R\). Define the sparse encoding
\[
  E_w(u)=\sum_{a\in I}u_aX^{w(a)}.
\]
Fix \(i,j\in I\). Suppose the selected exponent has exactly the expected
ordered-pair fibre:
\begin{equation}
\label{eq:fibre}
w(a)+w(b)=w(i)+w(j)
\quad\Longleftrightarrow\quad
(a=i\land b=j)\lor(a=j\land b=i).
\end{equation}
Then
\begin{equation}
\label{eq:coefficient}
[X^{w(i)+w(j)}]E_w(u)^2=
\begin{cases}
u_i^2,& i=j,\\
2u_i u_j,& i\ne j.
\end{cases}
\end{equation}

This needs only the fibre of one exponent. The weight map may have arbitrary
collisions elsewhere. It needs no positivity, cancellation, characteristic
zero, integral-domain assumption, or ordering on \(R\). In characteristic
two the off-diagonal coefficient may vanish; equation~\eqref{eq:coefficient}
still states the correct identity.

\subsection{Why the existing layering repeats work}

The global development in \code{QuadraticCoding} proves coefficient
extraction from the explicit cubic Sidon weights in \code{SidonPositions}.
The selective development in \code{SelectiveQuadraticCoding} works with an
arbitrary weight map and a local \code{Isolated} hypothesis. Both expand a
square into a double sum and prove essentially the same diagonal/off-diagonal
calculation \cite{quadratic,selective,sidon}.

The stronger global assumption is useful to establish
equation~\eqref{eq:fibre} for \emph{every} pair, but it is not needed by the
coefficient calculation itself. The new low-level module
\code{QuadraticFiber} states the calculation at its actual level of
generality. This avoids making the global module import the selective
module, which already depends on global notation and would create an
import cycle. Both instead use a common foundation containing only the
generic coefficient theorem.

\subsection{The checked interface}

The new theorem's hypothesis is local and exact:
\begin{lean}
theorem coeff_square_at_isolated
    [Fintype I] [DecidableEq I] [CommSemiring R]
    (w : I → Nat) (u : I → R) (i j : I)
    (hiso : ∀ a b, w a + w b = w i + w j ↔
      (a = i ∧ b = j) ∨ (a = j ∧ b = i)) :
    ((∑ a, Polynomial.monomial (w a) (u a)) ^ 2).coeff
      (w i + w j) = if i = j then u i ^ 2 else 2 * u i * u j
\end{lean}

The proof expands the finite product once, rewrites the coefficient test by
\code{hiso}, and separates the diagonal from the two distinct off-diagonal
terms. Keeping the diagonal case explicit is essential: when \(i=j\), the
two descriptions of the ordered pair refer to one term, not two.

The global diagonal and off-diagonal theorems specialize this with
\code{SidonPositions.weight_sum_eq_iff}. The global theorem on canonical
unordered pairs uses the same specialization directly. The selective
theorem becomes one application with its existing \code{Isolated} proof.
The public definitions of \code{encoded}, \code{pairWeight}, and the row
polynomials are preserved.

\subsection{Reuse the selective zero theorem for separated rows}

The same dependency correction has a useful second application. In
\code{MultiQuadraticCoding}, the theorem
\code{coeff_row_eq_zero_of_outside} repeats a nested
polynomial expansion to show that a coefficient outside a row's permitted
window is zero. The existing selective theorem
\code{coeff_at_other_target} already proves this conclusion whenever each
potential contributing exponent differs from the target
\cite{multi,selective}.

Specialize that theorem to the explicit Sidon weights and full pair
support. The client then retains only the arithmetic reason for
noncollision: the pair-weight bound, the bound on the sum of two weights,
and the lower bound on the row target. Its closing \code{omega} call checks
the interval arithmetic. There is no second sum/monomial proof in the
client.

This is part of the same direction rather than a separate claimed novelty:
general selective coefficient facts should sit below global special cases.

\subsection{Relation to excluded work and measured effect}

The earlier coefficient proposals concern upper envelopes, naturality,
support bounds, finite pushforward sums, and degree certificates. Those are
assumed to exist already. Here the new reusable statement is an \emph{exact}
coefficient formula from a two-point fibre, including the diagonal
multiplicity rule, and the proof graph is reorganized around it.

The integration patch adds one 32-line helper module and changes five
existing proof bodies across three modules. Aggregate source size is
reported with that helper included in Section~\ref{sec:validation}; blank
separators between declarations are preserved. The entire four-module
candidate was compiled, not merely the generic lemma in isolation.

The mathematical advantage outlasts a line-count reduction. A future
weight construction needs to prove the local fibre condition at its chosen
targets. It can then reuse the coefficient calculation even if unrelated
pair sums collide. A new global Sidon theorem is unnecessary for this
particular proof obligation.
